% !TeX root = ../../Book4.tex
% 纲要标题：### G.2 结合、交换、Lie 与 Poisson 运算
% * 状态：初稿已完成；新增范围尚未正式审核。
% #### G.2.1 Ass 与 Com 的运算和关系
% #### G.2.2 Lie 与 Poisson 的运算和关系
% #### G.2.3 代表性代数与低次核验

\section{结合、交换、Lie 与 Poisson 运算}
\label{b4:sec:G02}

同一个二元符号，施加不同的关系后会产生不同的代数理论。例如 \((ab)c=a(bc)\) 允许去掉乘积的括号，而 \(ab=ba\) 进一步允许忘掉输入次序。Operad 将这两种识别分别记录下来：前者改变复合树的形状，后者还改变对称群的作用。以下都在特征零域 \(k\) 上讨论，输入数至少为一；结合乘法与交换乘法均不预设单位元。Operad 的一元单位仍须作用为恒等映射，这与是否存在乘法单位是两个条件。

\subsection{Ass 与 Com 的运算和关系}
\label{b4:subsec:G02:01}

固定 \(n\ge1\)，考虑每个字母 \(x_1,\ldots,x_n\) 恰好出现一次、没有括号的字。它们有 \(n!\) 个；这些字先是形式符号，不是在某个给定代数中计算后得到的元素。

\begin{definition}\label{b4:def:ass-com-operads}
	设 \(k\) 为特征零域。对每个 \(n\ge1\)，以字母 \(x_1,\ldots,x_n\) 各使用一次的全部 \(n!\) 个字为基，所得向量空间记为 \(\operatorname{Ass}(n)\)。右 \(S_n\) 作用将旧标号 \(i\) 改成 \(\sigma^{-1}(i)\)。复合
	\[
		\gamma(w;w_1,\ldots,w_r)\in\operatorname{Ass}(n_1+\cdots+n_r)
	\]
	对 \(w\in\operatorname{Ass}(r)\) 与 \(w_i\in\operatorname{Ass}(n_i)\) 的基字由代入定义：把第 \(i\) 个输入块依次编号为 \(n_1+\cdots+n_{i-1}+1,\ldots,n_1+\cdots+n_i\)，再将 \(w\) 中的 \(x_i\) 换成相应块上的字 \(w_i\)。按多线性延拓，以字 \(x_1\) 为单位，所得 operad 称为\term[jiehe-operad]{结合 operad}[associative operad]。

	令 \(\operatorname{Com}(n)=k c_n\)，其中 \(S_n\) 作用平凡，并规定
	\[
		\gamma(c_r;c_{n_1},\ldots,c_{n_r})=c_{n_1+\cdots+n_r},
		\qquad \mathbf1=c_1.
	\]
	所得 operad 称为\term[jiaohuan-operad]{交换 operad}[commutative operad]。
\end{definition}
\symidx{\(\operatorname{Ass}\)：结合 operad}
\symidx{\(\operatorname{Com}\)：交换 operad}

例如
\[
	\gamma(x_2x_1;x_2x_1,x_1x_3x_2)=x_3x_5x_4x_2x_1.
\]
这里第一个输入块仍占标号 \(1,2\)，第二个输入块占标号 \(3,4,5\)；外层字 \(x_2x_1\) 使第二个块先出现。字的两次代入，无论先展开内层还是外层，最终都按照同一份嵌套次序排列全部字母，因此满足复合的结合律。代入单字 \(x_1\) 不改变原字。对输入重新编号再代入，与先代入再重新编号所得的字相同，故等变性也成立。\(\operatorname{Com}\) 的这些公理则归结为整数加法及标量乘法的结合律。

\ProseParagraphBreak
若将 \(\sigma\in S_n\) 对应于字
\[
	w_\sigma=x_{\sigma^{-1}(1)}\cdots x_{\sigma^{-1}(n)},
\]
则 \(w_\sigma\cdot\tau=w_{\sigma\tau}\)。因此 \(\operatorname{Ass}(n)\cong k[S_n]\) 是正则右模。这里取逆使记号与\subsecref{b4:subsec:G01:01}的右作用约定一致；字的形式仍更直接地显示运算中的输入次序。

\begin{proposition}\label{b4:prop:ass-com-algebras}
	设 \(k\) 为特征零域，\(A\) 为 \(k\) 向量空间。给定 \(A\) 上的 \(\operatorname{Ass}\) 代数结构，等价于给定一个不预设单位元的双线性结合乘法；给定 \(\operatorname{Com}\) 代数结构，等价于给定一个不预设单位元的双线性交换结合乘法。
\end{proposition}
\begin{proof}
	在 \(\operatorname{Ass}(2)\) 中记 \(\mu=x_1x_2\)。所有字均可由 \(\mu\)、一元单位及输入置换复合得到。两种三元复合满足
	\[
		\mu\circ_1\mu=x_1x_2x_3=\mu\circ_2\mu.
	\]
	所以 operad 态射 \(\theta:\operatorname{Ass}\to\operatorname{End}_A\) 必将 \(\mu\) 送到结合乘法 \(m\)，并把每个字送到按该字次序相乘的多线性映射。反过来，给定结合乘法，定义
	\[
		\theta(x_{i_1}\cdots x_{i_n})(a_1,\ldots,a_n)
		=a_{i_1}\cdots a_{i_n}.
	\]
	结合律使右边与加括号方式无关，也保证字的代入对应映射的复合。换输入次序与字母重编号一致，一元字作用为 \(\operatorname{id}_A\)，因而得到 operad 态射。两个构造互为逆。

	对 \(\operatorname{Com}\)，元素 \(c_2\) 在换位下不变，等变性于是要求 \(m(a,b)=m(b,a)\)；关系 \(c_2\circ_1c_2=c_2\circ_2c_2\) 要求结合律。反过来，交换结合乘法使 \(c_n\) 作用为 \(a_1\cdots a_n\)，并满足全部代入关系。这些构造都没有指定一个 \(0\) 元运算，所以没有产生乘法单位。
\end{proof}

这也给出两种 operad 的生成关系表示。在自由 operad 中取一个无对称限制的二元生成元，即令生成的右 \(S_2\) 模为 \(k[S_2]\)，再模掉 \(\mu\circ_1\mu-\mu\circ_2\mu\)，便得到 \(\operatorname{Ass}\)。确切地说，结合关系能将每棵二叉树改写为按同一次序相乘的右嵌套式；不同输入次序在字模型中线性无关，所以没有遗漏其他关系。

\ProseParagraphBreak
把每个字都送到 \(c_n\)，得到满态射
\[
	q:\operatorname{Ass}\longrightarrow\operatorname{Com}.
\]
其核是由 \(\mu-\mu\cdot(12)\) 生成的 operad 理想。事实上，结合关系允许将一对相邻输入先相乘，交换关系便能交换它们；相邻换位生成全部置换，因此商中所有 \(n\) 元字相等。它们的共同像 \(c_n\ne0\)，所以每个输入数的商恰为一维。\(q\) 所做的是增加交换关系；预复合 \(q\) 则把一个 \(\operatorname{Com}\) 代数看成 \(\operatorname{Ass}\) 代数。

\ProseParagraphBreak
一般自由代数公式\cref{b4:prop:free-operad-algebra}在这里恢复熟悉的对象。对向量空间 \(V\)，有
\[
	\begin{aligned}
		\operatorname{Ass}(V)&\cong\bigoplus_{n\ge1}V^{\otimes n}=T^+(V),\\
		\operatorname{Com}(V)&\cong\bigoplus_{n\ge1}\operatorname{Sym}^n(V)=\operatorname{Sym}^+(V).
	\end{aligned}
\]
第一个同构把 \(w_\sigma\otimes v\) 送到 \(\sigma v\)，利用的是 \(k[S_n]\otimes_{k[S_n]}V^{\otimes n}\cong V^{\otimes n}\)；第二个同构把张量模掉输入置换造成的差，因此得到对称幂。复合在第一式中是张量串接，在第二式中是对称乘法。上标 \(+\) 表示只取正次数部分：没有来自 \(V^{\otimes0}=k\) 的空字。

\subsection{Lie 与 Poisson 的运算和关系}
\label{b4:subsec:G02:02}

结合关系只含一种运算。Poisson 代数则要求一种乘法和一种括号同时存在，还要规定两者怎样相容。自由 operad 的商允许把这三类资料分别写清楚。

\begin{definition}\label{b4:def:lie-poisson-operads}
	设 \(k\) 为特征零域。在只由二元生成元 \(\lambda\) 生成的自由 operad 中，令 \(\lambda\cdot(12)=-\lambda\)，并模掉 Jacobi 元
	\[
		[[x_1,x_2],x_3]+[[x_2,x_3],x_1]+[[x_3,x_1],x_2],
	\]
	所生成的 operad 理想，其中 \([a,b]\) 表示生成元 \(\lambda\) 的代入。所得 operad 记为 \(\operatorname{Lie}\)，称为 \term[Lie-operad]{Lie operad}[Lie operad]。

	再取二元生成元 \(\mu,\lambda\)，其对称性分别为
	\[
		\mu\cdot(12)=\mu,\qquad \lambda\cdot(12)=-\lambda.
	\]
	在它们生成的自由 operad 中，模掉乘法的结合关系、括号的 Jacobi 关系及
	\begin{equation}\label{b4:eq:poisson-operad-leibniz}
		[x_1,x_2x_3]-[x_1,x_2]x_3-x_2[x_1,x_3]
	\end{equation}
	所生成的 operad 理想。这里并置表示 \(\mu\)，方括号表示 \(\lambda\)。所得 operad 记为 \(\operatorname{Pois}\)，称为 \term[Poisson-operad]{Poisson operad}[Poisson operad]。
\end{definition}
\symidx{\(\operatorname{Lie}\)：Lie operad}
\symidx{\(\operatorname{Pois}\)：Poisson operad}

生成元的反对称性已经写在生成右 \(S_2\) 模中；Jacobi 与 Leibniz 关系则位于输入数为三的部分。模掉 operad 理想还会模掉这些关系的全部代入和输入置换。因此商并非只保证三个独立字母满足恒等式，而是保证任意复合运算都保持这些恒等式。

\begin{proposition}\label{b4:prop:lie-poisson-algebras}
	设 \(k\) 为特征零域，\(A\) 为 \(k\) 向量空间。\(A\) 上的 \(\operatorname{Lie}\) 代数结构等价于 Lie 括号；\(A\) 上的 \(\operatorname{Pois}\) 代数结构等价于一个不预设单位元的交换结合乘法，以及对该乘法满足 Leibniz 法则的 Lie 括号。
\end{proposition}
\begin{proof}
	由\cref{b4:thm:free-operad}的泛性质，指定自由 operad 到 \(\operatorname{End}_A\) 的态射，只需指定生成元的像，并保持右 \(S_2\) 作用。\(\lambda\) 的像必须满足 \([a,b]=-[b,a]\)；特征零保证这与交替性等价。再由\cref{b4:prop:operad-quotient}，该态射通过 \(\operatorname{Lie}\) 当且仅当 Jacobi 元的像为零，也就是三个任意输入满足 Jacobi 恒等式。

	对 \(\operatorname{Pois}\)，\(\mu\) 的像同理是交换双线性乘法。三个关系族的像为零，依次就是结合律、Jacobi 恒等式和 Leibniz 法则。一旦这些生成关系为零，它们张成的全部 operad 理想也映为零：对零映射再作任意代入，所得仍为零。因此这些条件既必要又充分。
\end{proof}

\Cref{b4:def:poisson-algebra}对 Poisson 代数另要求乘法含幺；当前 operad 忘掉了这项要求。该处的每个 Poisson 代数都给出这里的 \(\operatorname{Pois}\) 代数。如果乘法原本已有单位，Leibniz 法则自动给出
\[
	\{a,1\}=\{a,1\cdot1\}=2\{a,1\},\qquad \{a,1\}=0;
\]
这说明括号与已有单位相容，却不能从中构造一个原来不存在的单位。

\ProseParagraphBreak
三个输入的关系都是两个二元运算的复合关系。这是它们在下一节成为二次 operad 的原因；此处的“二次”数的是树中的生成顶点，不是输入数。\(\operatorname{Pois}\) 的定义也说明，交换乘法与 Lie 括号单独成立还不够，必须另查\eqref{b4:eq:poisson-operad-leibniz}。

\subsection{代表性代数与低次核验}
\label{b4:subsec:G02:03}

若 \(A\) 为结合代数，则交换子 \([a,b]=ab-ba\) 给出 Lie 括号。反对称性直接成立，而
\[
	\begin{aligned}
		[[a,b],c]&=abc-bac-cab+cba,\\
		[[b,c],a]&=bca-cba-abc+acb,\\
		[[c,a],b]&=cab-acb-bca+bac
	\end{aligned}
\]
的和为零。这些等式在形式字中已经成立，所以有 operad 态射
\[
	\iota:\operatorname{Lie}\longrightarrow\operatorname{Ass},
	\qquad \lambda\longmapsto\mu-\mu\cdot(12).
\]
预复合 \(\iota\) 把结合代数变成其交换子 Lie 代数。箭头从记录 Lie 运算的 operad 指向记录结合运算的 operad，因为我们用结合乘法的表达式实现了 Lie 括号；所得代数结构的限制方向则相反。

\ProseParagraphBreak
任意 Lie 代数配上零乘法都是 \(\operatorname{Pois}\) 代数，因为 Leibniz 法则两边均为零；任意交换结合代数配上零括号也是 \(\operatorname{Pois}\) 代数。两种运算都不为零的例子是 \(k[q,p]\)，其中
\[
	\{f,g\}=\frac{\partial f}{\partial q}\frac{\partial g}{\partial p}
	-\frac{\partial f}{\partial p}\frac{\partial g}{\partial q}.
\]
\Cref{b4:prop:commuting-derivation-star-product}及\eqref{b4:eq:standard-symplectic-poisson}已验证它是 Poisson 括号；这里因此得到一个具体的 \(\operatorname{Pois}\) 代数。对这一已有例子，operad 记录的是普通乘法和括号及其恒等式，不包含星积的全部高阶系数。

\begin{proposition}\label{b4:prop:operad-low-arities}
	设 \(k\) 为特征零域。上述四个 operad 的低次维数为
	\[
		\begin{array}{c|rrrr}
			&\operatorname{Ass}&\operatorname{Com}&\operatorname{Lie}&\operatorname{Pois}\\\hline
			n=2&2&1&1&2\\
			n=3&6&1&2&6
		\end{array}
	\]
	其中 \(\operatorname{Lie}(3)\) 的一组基为
	\[
		\ell_1=[[x_1,x_2],x_3],\qquad
		\ell_2=[[x_2,x_3],x_1].
	\]
	\(\operatorname{Pois}(3)\) 的一组基为
	\[
		x_1x_2x_3,\quad [x_1,x_2]x_3,\quad [x_1,x_3]x_2,
		\quad [x_2,x_3]x_1,\quad \ell_1,\quad\ell_2.
	\]
\end{proposition}
\begin{proof}
	\(\operatorname{Ass}\) 与 \(\operatorname{Com}\) 的维数由定义给出。\(\operatorname{Lie}\) 的二元生成空间是一维，\(\operatorname{Pois}\) 的二元生成空间是对称及反对称的两份一维空间。三元关系及其代入不能降低输入数，所以没有二元生成元因此消失。

	反对称性将任意三元嵌套括号化为
	\[
		\ell_1,\quad \ell_2,\quad \ell_3=[[x_3,x_1],x_2]
	\]
	的线性组合；Jacobi 关系给出 \(\ell_3=-\ell_1-\ell_2\)。为证明剩余两个元素独立，将它们经 \(\iota\) 送入字模型。所得为
	\[
		\begin{aligned}
			\iota(\ell_1)&=x_1x_2x_3-x_2x_1x_3-x_3x_1x_2+x_3x_2x_1,\\
			\iota(\ell_2)&=x_2x_3x_1-x_3x_2x_1-x_1x_2x_3+x_1x_3x_2.
		\end{aligned}
	\]
	在 \(a\iota(\ell_1)+b\iota(\ell_2)\) 中，字 \(x_2x_1x_3\)、\(x_2x_3x_1\) 的系数分别为 \(-a,b\)，所以只能有平凡线性关系。这也证明 \(\iota\) 在输入数为三时为单射。

	对 \(\operatorname{Pois}(3)\)，按复合中括号的个数分成零个、一个、两个三部分。所有定义关系对此次数齐次，故商仍是这三部分的直和。无括号时，结合律与交换律给出单一生成元 \(x_1x_2x_3\)；它在通常的交换代数中可取非零值，因此不为零。有一个括号时，Leibniz 法则把括号在外的式子展开，再用乘法交换性与括号反对称性，得到所列三个“括号乘剩余输入”的生成元。在 \(k[q,p]\) 中分别代入
	\[
		(x_1,x_2,x_3)=(q,p,1),\quad(q,1,p),\quad(1,q,p),
	\]
	三次得到这三个生成元的值依次为 \((1,0,0)\)、\((0,1,0)\)、\((0,0,1)\)。因此它们线性无关。

	两个括号的部分由 \(\ell_1,\ell_2\) 生成。把自由结合代数的交换子 Lie 代数配上零乘法，就得到一个 \(\operatorname{Pois}\) 代数；在其中评价刚才的两个嵌套括号，仍得到两组线性无关的字。因此它们在 \(\operatorname{Pois}(3)\) 中也独立，合起来共有六个基元素。
\end{proof}

这六个三元 Poisson 运算分别对应三种分组：三个输入全部相乘，将其中两个先作括号再乘第三个，以及三个输入全放在嵌套括号中。Leibniz 关系负责把乘法从括号内部展开，Jacobi 关系负责消去嵌套括号之间的冗余。低次计算因而能直接显示各类关系的作用；\(\operatorname{Ass}(3)\) 与 \(\operatorname{Pois}(3)\) 虽然同为六维，它们的复合与对称群作用却并未由这个数字决定。
